\section{Conclusion and Future Work}

We introduced a method for diagnosing and improving causal abstractions by bucketing the input space according to interchange-intervention behavior. Rather than treating IIA as a single global summary, our method identifies subsets of inputs on which a proposed abstraction is nearly perfect. This turns causal abstraction into a more diagnostic and constructive framework
% : it reveals where a hypothesis works, where it fails, and what distinguishes the two cases
. Across fine-tuned and pretrained models of different sizes, and across alignment methods including DAS, MDAS, and full-vector patching, our experiments show that bucketing is a broadly applicable diagnostic tool for causal abstraction; in the toy logic task, it also supports iterative discovery of the high-level hypothesis itself.

Our approach has two main limitations. First, it depends on a prior task construal: it operates within a specified task and candidate abstraction, rather than performing unconstrained automatic interpretation. In this sense, it is better understood as a structured program search than as a general solution to mechanistic interpretability. Second, our experiments focus on relatively simple, mostly single-variable hypotheses. In more complex settings, interacting variables and accumulated approximation error may make it difficult to discover large buckets with high IIA. 
These limitations suggest two natural directions for future work. One is to scale the method to richer tasks, larger input spaces, and higher-dimensional or genuinely multi-variable abstractions. Another is to better understand how bucketing can support more automated hypothesis discovery and refinement in non-trivial settings---for example, by turning the structure of identified buckets into a systematic search procedure over candidate variables, decompositions, and compositions of partial hypotheses.

% We introduced a method for diagnosing and improving causal abstractions by bucketing the input space according to interchange-intervention behavior. Rather than treating interchange intervention accuracy as a single global summary, our method identifies subsets of inputs on which a proposed abstraction is nearly perfect. This shift turns causal abstraction into a more diagnostic and constructive framework: it not only measures whether a hypothesis works on average, but also reveals where it works, where it fails, and what distinguishes the two cases. In particular, across fine-tuned and pretrained models of different sizes, and across alignment methods including DAS, MDAS, and full-vector patching, our experiments show that the proposed bucketing procedure provides a broadly applicable diagnostic tool for causal abstraction; moreover, in the toy logic task, we demonstrate how our method supports iterative refinement of the high-level hypothesis itself.

% At the same time, our approach has some limitations. First, it depends on a prior task construal. The method does not perform unconstrained automatic interpretation of a model; rather, it operates within a specified task and a corresponding candidate abstraction. In this sense, it should be understood as a form of structured program search or hypothesis refinement, not as a general-purpose solution to mechanistic interpretability. Second, our experiments focus on relatively simple hypotheses and are largely single-variable in scope at each diagnosis step. While this regime is sufficient to demonstrate the core utility of bucketing, more complex hypotheses may involve interacting variables and compounding approximation errors, making the discovery of non-trivially large buckets with high IIA difficult.

% These limitations point directly to several promising directions for future work. A first priority is to scale the method to more realistic settings with richer tasks, larger input spaces, and higher-dimensional or genuinely multi-variable abstractions. A second direction is to better understand how bucketing can be used not only for diagnosis, but for automating non-trivial hypothesis discovery and refinement. In particular, an important open question is how to turn the signals revealed by the identified buckets into a systematic search procedure over candidate variables, decompositions, and compositions of partial hypotheses. We view this as a natural next step toward making causal-abstraction-style interpretability more scalable, more constructive, and less dependent on hand-specified hypotheses.